\chapter{Operational Graph Dynamics and Graph-Time Resonance}

The graph should not appear only as a static stiffness screen. In a
converter-rich power system, the graph is the spatial support on which power
flow, inertia, damping, controller memory, resonance, fragility, and design
act together. This chapter builds the missing dynamic graph layer. Its central
operator is
\begin{equation}
  T(s;z^\star)
  =
  s^2M+sD+B W(z^\star)B^T+\Pi(s;z^\star),
  \label{eq:graph_dynamic_master_operator}
\end{equation}
where \(B\) is an oriented incidence matrix, \(W(z^\star)\) contains
incremental line stiffnesses at the AC equilibrium \(z^\star\), \(M\) is a
nodal inertia measure, \(D\) is a nodal or device-level dissipative field, and
\(\Pi(s;z^\star)\) is the condensed controller self-energy. This turns the
thesis into a theory of damped, controlled waves on electrical graphs.

\thesisbox{thesisblue}{Graph-dynamic thesis}{
Dynamic stability is a joint property of topology, operating point, spatial
inertia and damping fields, and controller memory. A resonance occurs when
temporal proximity, spatial alignment, port visibility, low damping, and
resolvent amplification coincide.}

\section{Power-Flow Geometry of the Operational Graph}

Consider a connected lossless graph with \(n\) buses, \(m\) lines, and an
arbitrarily oriented incidence matrix \(B\in\R^{n\times m}\). Let
\begin{equation}
  \eta=B^T\theta
  \label{eq:graph_eta}
\end{equation}
be the vector of line angle differences. For line \(e=(i,j)\), define
\begin{equation}
  \kappa_e=b_eV_iV_j>0.
  \label{eq:graph_kappa}
\end{equation}
The oriented active line flow is
\begin{equation}
  f_e(\theta,V)=\kappa_e\sin\eta_e,
  \label{eq:graph_line_flow}
\end{equation}
and the nodal active-power injection map is
\begin{equation}
  p(\theta,V)=B\operatorname{diag}(\kappa)\sin(B^T\theta).
  \label{eq:graph_power_flow_compact}
\end{equation}
At an equilibrium \(z^\star=(\theta^\star,V^\star,\ldots)\),
\begin{equation}
  f_e^\star=\kappa_e^\star\sin\eta_e^\star,
  \qquad
  w_e^\star=\kappa_e^\star\cos\eta_e^\star.
  \label{eq:graph_flow_stiffness_pair}
\end{equation}
Thus the operating-point stiffness is
\begin{equation}
  L^\star=B W^\star B^T,\qquad
  W^\star=\operatorname{diag}(w_e^\star).
  \label{eq:graph_operational_laplacian}
\end{equation}
If \(|\eta_e^\star|<\pi/2\) for all lines, then \(W^\star\succ0\). For a
connected graph, \(L^\star\succeq0\) and
\(\ker L^\star=\operatorname{span}\{\ones\}\). The uniform angle direction is
removed by projection onto \(\ones^\perp\).

The flow-stiffness pair in \eqref{eq:graph_flow_stiffness_pair} is the first
dynamic graph fact:
\begin{equation}
  \frac{\partial w_e}{\partial \eta_e}
  =
  -\kappa_e\sin\eta_e
  =
  -f_e.
  \label{eq:graph_stiffness_angle_derivative}
\end{equation}
A heavily loaded line can lose synchronizing stiffness rapidly under an
additional angular perturbation. This is the direct bridge from power flow to
modal fragility.

\section{Total Derivative of the Operational Laplacian}

Let \(\rho\) be a planning parameter that may change injections,
susceptances, voltages, topology, or the operating point. The reduced
equilibrium equation is
\begin{equation}
  F(\theta,\rho)
  =
  B\operatorname{diag}(\kappa(\rho))\sin(B^T\theta)-p(\rho)=0.
  \label{eq:graph_equilibrium_F}
\end{equation}
Assume the reduced Laplacian \(L\) is invertible on \(\ones^\perp\). Then
\begin{equation}
  \theta_\rho
  =
  L^\dagger
  \left[
  p_\rho-B\operatorname{diag}(\sin\eta)\kappa_\rho
  \right],
  \qquad
  \eta_\rho=B^T\theta_\rho.
  \label{eq:graph_theta_rho}
\end{equation}
Since \(w=\kappa\odot\cos\eta\),
\begin{equation}
  w_\rho
  =
  \cos\eta\odot\kappa_\rho
  -
  f\odot\eta_\rho,
  \label{eq:graph_w_rho}
\end{equation}
and
\begin{equation}
  L_\rho=B\operatorname{diag}(w_\rho)B^T.
  \label{eq:graph_L_rho}
\end{equation}

Equation~\eqref{eq:graph_w_rho} separates the direct line-parameter channel
from the operating-point channel:
\begin{equation}
  L_\rho
  =
  B\operatorname{diag}(\cos\eta\odot\kappa_\rho)B^T
  -
  B\operatorname{diag}(f\odot\eta_\rho)B^T.
  \label{eq:graph_L_rho_split}
\end{equation}
The second term is often the one missed by frozen-graph screens. A physical
line change modifies flows and angles; those changes modify many incremental
line stiffnesses, not only the line being touched.

\section{Mode Sensitivity to Power-Flow Geometry}

For generalized graph modes
\begin{equation}
  L\phi_k=\nu_kM\phi_k,
  \qquad
  \phi_k^TM\phi_\ell=\delta_{k\ell},
  \label{eq:graph_generalized_modes}
\end{equation}
a simple eigenvalue satisfies
\begin{equation}
  (\nu_k)_\rho
  =
  \phi_k^T(L_\rho-\nu_kM_\rho)\phi_k.
  \label{eq:graph_nu_rho_general}
\end{equation}
Let \(a_e\) be the incidence vector of line \(e\), and define the modal line
deformation
\begin{equation}
  g_{e,k}=a_e^T\phi_k.
  \label{eq:graph_modal_line_deformation}
\end{equation}
Then
\begin{equation}
  (\nu_k)_\rho
  =
  \sum_{e=1}^m (w_e)_\rho g_{e,k}^2
  -
  \nu_k\sum_{i=1}^n (M_i)_\rho \phi_k(i)^2.
  \label{eq:graph_nu_rho_edges}
\end{equation}

If only injections change while \(\kappa\) and \(M\) are fixed, then
\(\theta_\rho=L^\dagger p_\rho\), and
\begin{equation}
  (\nu_k)_\rho
  =
  -
  \sum_e f_e g_{e,k}^2 a_e^TL^\dagger p_\rho.
  \label{eq:graph_injection_nu_rho}
\end{equation}
Equivalently,
\begin{equation}
  \nabla_p\nu_k
  =
  -L^\dagger B\operatorname{diag}(f)c_k,
  \qquad
  c_k=(g_{1,k}^2,\ldots,g_{m,k}^2)^T.
  \label{eq:graph_grad_p_nu}
\end{equation}
The flow-weighted modal centrality
\begin{equation}
  \chi_{e,k}^{\rm flow}
  =
  |f_e|\,|g_{e,k}|^2
  \label{eq:graph_flow_modal_centrality}
\end{equation}
flags lines that simultaneously carry high flow and sit across a large jump of
the critical mode. This is more physical than degree, betweenness, or
\(\partial\lambda_2/\partial b_e\) alone.

\section{Cycle and Cut Spaces}

The DC flow induced by \(p\in\ones^\perp\) is
\begin{equation}
  f^\star=W B^T L^\dagger p.
  \label{eq:graph_dc_flow}
\end{equation}
It is the unique solution of
\begin{equation}
  \min_f\ \frac12 f^TW^{-1}f
  \quad
  \text{subject to}\quad Bf=p.
  \label{eq:graph_min_energy_flow}
\end{equation}
The Lagrange multiplier proof is immediate: stationarity gives
\(f=WB^T\lambda\), and the constraint gives \(L\lambda=p\).

If \(C\) spans the cycle space, \(BC=0\), every feasible flow can be written
as \(f=f_0+Cq\). Changing a line weight changes the metric \(W^{-1}\), modifies
cycle redistribution, changes the operating point, and then changes \(L\).
This is the formal bridge between topology control, Braess-like effects, and
modal dynamics.

For a cut \(S\subset\mathcal V\), define
\begin{equation}
  m_S=\sum_{i\in S}M_i,\qquad
  m_{S^c}=\sum_{i\notin S}M_i,\qquad
  w(\partial S)=\sum_{e\in\partial S}w_e.
  \label{eq:graph_cut_quantities}
\end{equation}
The first nonzero generalized eigenvalue satisfies
\begin{equation}
  \nu_2
  \le
  w(\partial S)
  \left(
  \frac1{m_S}+\frac1{m_{S^c}}
  \right)
  \label{eq:graph_inertia_cut_bound}
\end{equation}
for every nontrivial cut. The proof is the Rayleigh quotient evaluated on the
two-area vector that is constant on \(S\) and \(S^c\) and has zero
center-of-inertia mean.

\section{Graph Roughness of Damping-to-Inertia Fields}

Assume \(M=\diag(M_i)\) and \(D=\diag(D_i)\), and define the nodal field
\begin{equation}
  \gamma_i=\frac{D_i}{M_i}.
  \label{eq:graph_gamma_field}
\end{equation}
Then
\begin{equation}
  \Dt=M^{-1/2}DM^{-1/2}=\diag(\gamma_i).
  \label{eq:graph_Dt_gamma}
\end{equation}
For \(i\ne j\),
\begin{equation}
  [\Lt,\Dt]_{ij}
  =
  \Lt_{ij}(\gamma_j-\gamma_i).
  \label{eq:graph_commutator_entry}
\end{equation}
Therefore
\begin{equation}
  \norm{[\Lt,\Dt]}_F^2
  =
  2\sum_{e=(i,j)}
  \frac{w_e^2}{M_iM_j}
  (\gamma_i-\gamma_j)^2.
  \label{eq:graph_roughness_identity}
\end{equation}
This is a Dirichlet energy over an auxiliary graph with weights
\(\widehat w_e=w_e^2/(M_iM_j)\). Within the diagonal nodal damping class, and
for a connected graph,
\begin{equation}
  [\Lt,\Dt]=0
  \quad\Longleftrightarrow\quad
  \gamma_i=\gamma\ \text{for all }i
  \quad\Longleftrightarrow\quad
  D=\gamma M.
  \label{eq:graph_diagonal_commuting_iff}
\end{equation}
Thus non-proportionality is spatial roughness of the damping-to-inertia field.

Let \(\Lt Q=Q\Lambda\) and \(\Gamma=Q^T\Dt Q\). Since
\begin{equation}
  Q^T[\Lt,\Dt]Q=[\Lambda,\Gamma],
  \label{eq:graph_commutator_modal}
\end{equation}
for \(k\ne\ell\),
\begin{equation}
  (\nu_k-\nu_\ell)\Gamma_{k\ell}
  =
  q_k^T[\Lt,\Dt]q_\ell.
  \label{eq:graph_modal_coupling_identity}
\end{equation}
Therefore
\begin{equation}
  |\Gamma_{k\ell}|
  \le
  \frac{\norm{[\Lt,\Dt]}_2}{|\nu_k-\nu_\ell|}
  \le
  \frac{
  \left(
  2\sum_{e=(i,j)}
  \frac{w_e^2}{M_iM_j}(\gamma_i-\gamma_j)^2
  \right)^{1/2}}
  {|\nu_k-\nu_\ell|}.
  \label{eq:graph_modal_mixing_bound}
\end{equation}
Large graph roughness and small spectral separation are the two ingredients of
strong modal mixing. The design reading is direct: smooth the field
\(D_i/M_i\), avoid abrupt local damping-to-inertia jumps across high-stiffness
edges, or escalate from a graph-mode damping screen to the full QEP whenever
\eqref{eq:graph_modal_mixing_bound} is not small in the protected mode band.

\section{Dynamic Roughness of Local Control Fields}

If the damping-channel self-energy is approximately local,
\begin{equation}
  \Sigma(s)=\diag(\sigma_i(s)),
  \qquad
  \gamma_i(s)=\frac{\sigma_i(s)}{M_i},
  \label{eq:graph_dynamic_gamma}
\end{equation}
then, frequency by frequency,
\begin{equation}
  \norm{[\Lt,\widetilde\Sigma(s)]}_F^2
  =
  2\sum_{e=(i,j)}
  \frac{w_e^2}{M_iM_j}
  |\gamma_i(s)-\gamma_j(s)|^2.
  \label{eq:graph_dynamic_roughness}
\end{equation}
The band function
\begin{equation}
  \mathcal E_\gamma(\omega)
  =
  \sum_{e=(i,j)}
  \frac{w_e^2}{M_iM_j}
  |\gamma_i(j\omega)-\gamma_j(j\omega)|^2
  \label{eq:graph_dynamic_roughness_band}
\end{equation}
identifies frequencies where the spatial distribution of control is
incompatible with the electrical topology. If \(\Pi(s)\) is dense, it should
be decomposed into edge-like, shunt-like, and irreducibly nonlocal parts:
\begin{equation}
  \Pi(j\omega)
  =
  B K_\Pi(j\omega)B^T
  +
  \diag(k_\Pi(j\omega))
  +
  \Pi_{\rm nonlocal}(j\omega).
  \label{eq:graph_pi_edge_shunt_nonlocal}
\end{equation}
This is the dynamic extension of the line-shunt decomposition used later for
static perturbations.

\section{Joint Graph-Time Frequency Response}

Normalize the retained operator:
\begin{equation}
  \widetilde T(s)=M^{-1/2}T(s)M^{-1/2}.
  \label{eq:graph_normalized_T}
\end{equation}
Let \(\Lt=Q\Lambda Q^T\). In graph-frequency coordinates,
\begin{equation}
  T_G(s)
  =
  Q^T\widetilde T(s)Q
  =
  s^2I+s\Gamma+\Lambda+\Pi_G(s),
  \label{eq:graph_time_operator}
\end{equation}
where \(\Pi_G(s)=Q^TM^{-1/2}\Pi(s)M^{-1/2}Q\). The graph-time transfer
function is
\begin{equation}
  H_G(\omega)=T_G(j\omega)^{-1}.
  \label{eq:graph_time_transfer}
\end{equation}
The index \(k\) is a spatial graph frequency, while \(\omega\) is temporal
frequency. \(H_{G,kk}(\omega)\) measures amplification of spatial mode \(k\),
and \(H_{G,k\ell}(\omega)\) measures conversion from spatial mode \(\ell\) to
spatial mode \(k\).

Let
\begin{equation}
  d_k(\omega)
  =
  \nu_k-\omega^2
  +j\omega\Gamma_{kk}
  +[\Pi_G(j\omega)]_{kk},
  \label{eq:graph_dk}
\end{equation}
and write
\begin{equation}
  T_G(j\omega)=D_0(\omega)+E(\omega),
  \qquad
  D_0(\omega)=\diag(d_1,\ldots,d_n).
  \label{eq:graph_TG_split}
\end{equation}
If \(\alpha(\omega)=\norm{D_0(\omega)^{-1}E(\omega)}_2<1\), then
\begin{equation}
  \norm{H_G(\omega)-D_0(\omega)^{-1}}_2
  \le
  \frac{\norm{D_0^{-1}}_2^2\norm{E}_2}{1-\alpha(\omega)}.
  \label{eq:graph_cross_resonance_bound}
\end{equation}
To first order,
\begin{equation}
  [H_G(\omega)]_{k\ell}
  \approx
  -\frac{E_{k\ell}(\omega)}{d_k(\omega)d_\ell(\omega)},
  \qquad k\ne\ell.
  \label{eq:graph_cross_resonance_first_order}
\end{equation}
A cross-modal resonance therefore requires temporal amplification in both
modes and a strong off-diagonal channel from damping, self-energy, nonnormal
AC stiffness, or nonlocal control.

\section{Resonance Taxonomy}

The graph-time response separates six mechanisms:
\begin{enumerate}
\item \textbf{Natural resonance:} \(\det T(s)=0\).
\item \textbf{Forced resonance:} \(j\Omega\) lies near a visible pole and the
input shape has strong projection on the corresponding left eigenvector.
\item \textbf{Cross-modal graph resonance:} \(|H_{G,k\ell}(\omega)|\) is large
for \(k\ne\ell\).
\item \textbf{Network-control resonance:} a pole of \(\Pi(s)\) is close to a
network-family pole and has a visible residue.
\item \textbf{Parametric topology resonance:} a time-varying line weight
modulates \(L(t)=L_0+\varepsilon L_1\cos\Omega t\).
\item \textbf{Internal nonlinear resonance:} finite-amplitude dynamics satisfy
\(m\omega_k\approx n\omega_\ell\) and require normal forms, continuation, or
harmonic balance.
\end{enumerate}

For parametric line modulation,
\begin{equation}
  L_1=B\widehat W B^T,\qquad
  c_{k\ell}=\phi_k^TL_1\phi_\ell
  =
  \sum_e \widehat w_e(a_e^T\phi_k)(a_e^T\phi_\ell).
  \label{eq:graph_parametric_coupling}
\end{equation}
The coefficient \(c_{k\ell}\) is the exact spatial bridge from a modulated
edge to the two interacting graph modes.

\section{Dynamic Effective Resistance and Low-Rank Updates}

Let
\begin{equation}
  G(s)=T(s)^{-1}
  \label{eq:graph_green_function}
\end{equation}
on the projected nonzero-angle subspace. For line \(e\), define the dynamic
effective resistance
\begin{equation}
  \mathcal R_e(s)=a_e^TG(s)a_e.
  \label{eq:graph_dynamic_effective_resistance}
\end{equation}
If \(s=0\), \(\Pi(0)=0\), and \(T(0)=L\), this reduces to the static effective
resistance \(a_e^TL^\dagger a_e\). For \(s\ne0\), it contains topology,
operating point, inertia, damping, controller poles, memory, and resolvent
amplification.

For a frozen operating point and a rank-one stiffness update,
\begin{equation}
  T_\Delta(s)=T(s)+\Delta w_ea_ea_e^T,
  \label{eq:graph_rank_one_update}
\end{equation}
the determinant lemma gives
\begin{equation}
  \det T_\Delta(s)
  =
  \det T(s)\left[1+\Delta w_e\mathcal R_e(s)\right].
  \label{eq:graph_det_rank_one}
\end{equation}
New or shifted poles caused by this finite frozen update satisfy
\begin{equation}
  1+\Delta w_e\mathcal R_e(s)=0.
  \label{eq:graph_pole_update_equation}
\end{equation}
The corresponding Green function update is
\begin{equation}
  G_\Delta(s)
  =
  G(s)
  -
  \frac{\Delta w_e\,G(s)a_ea_e^TG(s)}
  {1+\Delta w_e\mathcal R_e(s)}.
  \label{eq:graph_green_update}
\end{equation}

For multiple candidate lines with incidence matrix \(A_C\),
\begin{equation}
  K_C(s)=A_C^TG(s)A_C,
  \label{eq:graph_candidate_kernel}
\end{equation}
and
\begin{equation}
  \det T_\Delta(s)
  =
  \det T(s)\det\!\left[I+\Delta W K_C(s)\right].
  \label{eq:graph_multi_det_update}
\end{equation}
This reduces a large spectral update to a matrix whose dimension is the number
of candidate actions. It is exact for frozen low-rank stiffness updates and a
predictor for physical line changes that also move the operating point.

\section{Dynamic Edge Centrality}

For an input shape \(b\), an output shape \(c\), and
\begin{equation}
  h(s)=c^TG(s)b,
  \label{eq:graph_transfer_scalar}
\end{equation}
the frozen line-weight derivative is
\begin{equation}
  \frac{\partial h(s)}{\partial w_e}
  =
  -
  \left[c^TG(s)a_e\right]
  \left[a_e^TG(s)b\right].
  \label{eq:graph_dynamic_edge_derivative}
\end{equation}
Thus
\begin{equation}
  \mathcal C_e(\omega;b,c)
  =
  |c^TG(j\omega)a_e|\,
  |a_e^TG(j\omega)b|
  \label{eq:graph_dynamic_edge_centrality}
\end{equation}
is a frequency-specific dynamic centrality. It identifies a line that is
reachable from the disturbance, transmits energy toward the protected output,
and matters in the resonant band. It need not have high static degree,
betweenness, or Fiedler sensitivity.

\section{Structured Fragility Radii}

There is no single maximum perturbation. At least four different radii are
needed.

For topology or parameter changes, define a protected spectral region
\(\mathcal S\). With candidate line matrix \(A_C\), the structured topological
radius is
\begin{equation}
  r_{\rm top}
  =
  \inf\left\{
  \norm{\Delta W}:
  \exists s\in\partial\mathcal S,\ 
  \det[T(s)+A_C\Delta W A_C^T]=0
  \right\}.
  \label{eq:graph_topological_radius}
\end{equation}
Using \eqref{eq:graph_multi_det_update}, the reduced condition is
\begin{equation}
  \det[I+\Delta W K_C(s)]=0.
  \label{eq:graph_reduced_fragility_condition}
\end{equation}
With a declared physical uncertainty structure, this becomes the structured
radius
\begin{equation}
  r_{\rm top}
  =
  \left[
  \sup_{s\in\partial\mathcal S}
  \mu_\Delta(K_C(s))
  \right]^{-1}.
  \label{eq:graph_mu_radius}
\end{equation}

For a single frozen line update, the complex unstructured line radius along a
protected boundary is especially transparent:
\begin{equation}
  r_{e}^{\C}
  =
  \left[
  \sup_{s\in\partial\mathcal S}
  |\mathcal R_e(s)|
  \right]^{-1}.
  \label{eq:graph_complex_line_radius}
\end{equation}
This is the smallest complex scalar \(\Delta w_e\) that can satisfy
\(1+\Delta w_e\mathcal R_e(s)=0\) for some boundary point. A physical line
weight is real, so the real frozen-line radius is
\begin{equation}
  r_{e}^{\R}
  =
  \min_{\substack{s\in\partial\mathcal S\\
  \operatorname{Im}\mathcal R_e(s)=0}}
  \frac{1}{|\operatorname{Re}\mathcal R_e(s)|},
  \label{eq:graph_real_line_radius}
\end{equation}
with the convention \(r_e^{\R}=\infty\) if the feasibility set is empty. Thus
dynamic effective resistance is not only a ranking score; it gives an exact
single-line fragility radius for the declared spectral boundary.

For a forced input in a stable linear model, increasing amplitude does not move
poles. It can still violate an operational limit:
\begin{equation}
  |\widehat u|_{\max}^{\Omega}
  =
  \frac{y_{\max}}
  {\sup_{\omega\in\Omega}
  \sigma_{\max}\!\left(C_yG(j\omega)B_u\right)}.
  \label{eq:graph_forced_amplitude_radius}
\end{equation}
For nonlinear loss of stability, the relevant object is a region-of-attraction
radius,
\begin{equation}
  r_{\rm ROA}
  =
  \inf_{x\in\partial\mathcal B(x^\star)}
  \norm{W_x(x-x^\star)}.
  \label{eq:graph_roa_radius}
\end{equation}
The first is a parametric fragility radius, the second is an input-output
exposure bound, and the third is a nonlinear stability quantity.

\section{Stability Recovery Design}

The graph-dynamic design problem combines line, damping, inertia, and
controller actions:
\begin{equation}
  u=
  \begin{bmatrix}
  \Delta w & \Delta D & \Delta M & \Delta\kappa
  \end{bmatrix}^T.
  \label{eq:graph_design_vector}
\end{equation}
A unified recovery problem is
\begin{align}
  \min_u\quad&
  c_w^T|\Delta w|
  +c_D^T|\Delta D|
  +c_M^T|\Delta M|
  +c_\kappa^T|\Delta\kappa|
  \label{eq:graph_recovery_problem}\\
  \text{subject to}\quad&
  F(z^\star,u)=0,\nonumber\\
  &|f_e(z^\star,u)|\le \overline f_e,\nonumber\\
  &|B^T\theta^\star|\le \eta_{\max},\nonumber\\
  &\Real s_k(u)\le-\alpha_\star,\qquad
    \zeta_k(u)\ge\zeta_\star,\nonumber\\
  &\sup_{\omega\in\Omega}
    \sigma_{\max}\!\left(C_yT(j\omega,u)^{-1}B_u\right)
    \le g_\star,\nonumber\\
  &r_{\rm top}(u)\ge r_\star.\nonumber
\end{align}
Sequential versions should use pole sensitivities and resolvent-gradient
steps only inside a trust region, followed by continuation of the operating
point and the protected poles.

\section{AC, Directed, and Matrix-Weighted Extensions}

The scalar lossless Laplacian is the symmetric floor, not the universal model.
For the AC Jacobian,
\begin{equation}
  \begin{bmatrix}
  \Delta P\\
  \Delta Q
  \end{bmatrix}
  =
  \begin{bmatrix}
  J_{P\theta}&J_{PV}\\
  J_{Q\theta}&J_{QV}
  \end{bmatrix}
  \begin{bmatrix}
  \Delta\theta\\
  \Delta V
  \end{bmatrix},
  \label{eq:graph_ac_jacobian}
\end{equation}
eliminating voltage increments under \(\Delta Q=0\) gives
\begin{equation}
  L_{\rm AC}
  =
  J_{P\theta}-J_{PV}J_{QV}^{-1}J_{Q\theta}.
  \label{eq:graph_ac_condensed_L}
\end{equation}
In general \(L_{\rm AC}\ne L_{\rm AC}^T\). Its symmetric part is
stiffness-like; its skew part is circulatory and nonconservative. This
nonnormal layer requires left-right eigenvectors and resolvent gains.

For multivariable ports, assign each bus an \(r\)-dimensional vector and set
\(\mathcal B=B\otimes I_r\). With line admittances
\(Y_e(s)\in\C^{r\times r}\),
\begin{equation}
  \mathcal L(s)
  =
  \mathcal B
  \operatorname{blkdiag}(Y_1(s),\ldots,Y_m(s))
  \mathcal B^T.
  \label{eq:graph_matrix_laplacian}
\end{equation}
For a matrix update on line \(e\), with \(\mathcal A_e=a_e\otimes I_r\),
\begin{equation}
  \mathcal R_e^{(r)}(s)
  =
  \mathcal A_e^T\mathcal T(s)^{-1}\mathcal A_e,
  \label{eq:graph_matrix_effective_resistance}
\end{equation}
and
\begin{equation}
  \det\mathcal T_\Delta(s)
  =
  \det\mathcal T(s)
  \det\!\left[I_r+\Delta Y_e(s)\mathcal R_e^{(r)}(s)\right].
  \label{eq:graph_matrix_det_update}
\end{equation}
This is the natural extension to \(dq\) coupling, voltage-frequency channels,
frequency-dependent line models, directed effects, and multivariable converter
ports.

\section{Paper Extraction}

This chapter supports a focused Transactions-style paper:
\begin{quote}
\itshape
Graph-Time Resonance, Spectral Effective Resistance, and Structured Topology
Fragility in Converter-Rich Power Systems.
\end{quote}
The broad dissertation contribution is operator-valued dynamic grid strength.
The narrower article should be sold as controller-resolvent-aware spectral
graph planning. Its minimum publishable package is:
\begin{enumerate}
\item graph roughness of \(D_i/M_i\) and the modal-mixing bound
\eqref{eq:graph_modal_mixing_bound};
\item graph-time transfer \(H_G(\omega)\) and cross-resonance atlas;
\item finite dynamic-effective-resistance updates
\(1+\Delta w_e\mathcal R_e(s)=0\);
\item structured topology-fragility radii
\eqref{eq:graph_complex_line_radius}--\eqref{eq:graph_real_line_radius};
\item a controller-induced reinforcement-reversal gate, using the Schur
self-energy machinery developed later in the thesis.
\end{enumerate}
The novelty claim should not be that Laplacians or effective resistance are
new, that Schur complements are new, or that adding a line can never hurt in
existing literature. The defensible claim is narrower and stronger:
power-flow-induced stiffness, damping-to-inertia roughness, graph-time modal
conversion, finite low-rank topology updates, rational controller self-energy,
and structured stability radii are projections of one operator
\(T(s;z^\star)\), and that operator determines when a topological improvement
really improves a protected dynamic margin.
